% Worked Examples: Concept 2.4.1 - Reachability
\documentclass[11pt,a4paper]{article}
\usepackage{worked-examples-template}

\title{\textbf{Worked Examples}\\[0.5em]
\large Concept 2.4.1: Reachability\\[0.3em]
\normalsize \coursesubtitle{} -- \coursetitle}
\author{Assoc.\ Prof.\ William Robertson \and Dr Sean McGowan}
\date{\today}

\begin{document}

\maketitle
\tableofcontents
\newpage

\section{Concept 2.4.1: Reachability}

\subsection{Example 1: Testing Reachability of a Second-Order System}

\paragraph{Problem:} Consider a simple mechanical system described by the state-space model:
\begin{align}
\frac{d}{dt}\begin{bmatrix} x_1 \\ x_2 \end{bmatrix} = \begin{bmatrix} 0 & 1 \\ -2 & -3 \end{bmatrix}\begin{bmatrix} x_1 \\ x_2 \end{bmatrix} + \begin{bmatrix} 0 \\ 1 \end{bmatrix}u
\end{align}
where $x_1$ is position and $x_2$ is velocity. Determine if this system is reachable.

\paragraph{Solution:} 

To test reachability, we compute the reachability matrix $W_r$ and check if it is full rank (invertible).

For a system of dimension $n=2$, the reachability matrix is:
\begin{align}
W_r = \begin{bmatrix} B & AB \end{bmatrix}
\end{align}

Given:
\begin{align}
A = \begin{bmatrix} 0 & 1 \\ -2 & -3 \end{bmatrix}, \quad B = \begin{bmatrix} 0 \\ 1 \end{bmatrix}
\end{align}

First, compute $AB$:
\begin{align}
AB = \begin{bmatrix} 0 & 1 \\ -2 & -3 \end{bmatrix}\begin{bmatrix} 0 \\ 1 \end{bmatrix} = \begin{bmatrix} 1 \\ -3 \end{bmatrix}
\end{align}

Therefore, the reachability matrix is:
\begin{align}
W_r = \begin{bmatrix} 0 & 1 \\ 1 & -3 \end{bmatrix}
\end{align}

To check if $W_r$ is invertible, compute the determinant:
\begin{align}
\det(W_r) = (0)(-3) - (1)(1) = -1 \neq 0
\end{align}

Since $\det(W_r) \neq 0$, the reachability matrix is full rank and invertible.

\textbf{Conclusion:} The system is reachable. This means we can steer the system from any initial state to any desired final state in finite time using an appropriate control input $u(t)$.
\end{document}
